\chapter{From Source to Release: CI/CD Through This Textbook}
\label{app:cicd}

This appendix follows one change through the automation that builds and
publishes this book. The example is small enough to inspect completely, but it
contains the same boundaries found in a scientific service or machine-learning
product: reviewed source, clean builds, immutable artifacts, named
environments, scoped authority, post-deployment verification, and recovery.

No prior knowledge of CI/CD is assumed. The objective is not to memorize a
particular GitHub Actions syntax. It is to learn which claim each stage makes,
which evidence supports that claim, and which authority allows the system to
change public state.

After working through this appendix, you should be able to:

\begin{itemize}
  \item distinguish continuous integration, continuous delivery, and
  continuous deployment;
  \item distinguish source, a build, an artifact, a release, a deployment, and
  an environment;
  \item trace a pull request, a merge to \code{main}, and a version tag through
  this repository's workflows;
  \item explain why a later stage should promote the artifact that an earlier
  stage verified rather than rebuild it;
  \item read a deployment manifest and verify an artifact with a checksum;
  \item identify the permissions and human decisions that control publication;
  and
  \item design a rollback or roll-forward response that preserves evidence.
\end{itemize}

\section{A public system has more state than its source}

A source file is not a deployed system. The LaTeX sentence being read now is
source. The compiled PDF is an artifact. A named collection containing that
PDF, a Python wheel, and a source archive is a release bundle. A copy of the PDF
served from the public course site is a deployment. These objects can agree or
disagree. A disciplined pipeline makes their relationships explicit.

\begingroup\small
\begin{longtable}{@{}
  >{\raggedright\arraybackslash}p{0.19\textwidth}
  >{\raggedright\arraybackslash}p{0.35\textwidth}
  >{\raggedright\arraybackslash}p{0.36\textwidth}@{}}
\toprule
\textbf{Object} & \textbf{Meaning} & \textbf{Identity in this repository} \\
\midrule
\endhead
Source revision & One immutable state of version-controlled inputs & Complete
Git commit identifier \\
Build & Process that transforms source into outputs & A particular workflow
job and its logs \\
Artifact & Bounded output intended to move between stages & Static-site
archive, PDF, wheel, or source distribution \\
Environment & Named destination with configuration and policy &
\code{github-pages} or \code{course-release} \\
Deployment & Placement of an artifact into an environment & One GitHub Pages
deployment record \\
Release & Named statement that artifacts represent a version & GitHub Release
attached to \code{course-vX.Y.Z} \\
\bottomrule
\end{longtable}
\endgroup

The word \emph{environment} is overloaded. A Python virtual environment
isolates an interpreter and installed packages. A deployment environment names
a destination such as staging or production and may impose approval,
configuration, and secret policies. Context must identify which meaning is
intended.

\begin{professionalbox}
A branch name such as \code{main} is movable. It is convenient for navigation,
but it is not sufficient artifact identity. A defensible record binds public
state to an immutable commit identifier and, when relevant, an artifact digest
and release tag.
\end{professionalbox}

\section{The three terms in CI/CD}

\term{Continuous integration} (CI) automatically tests an integrated source
revision in a clean environment. A pull request is a candidate integration. CI
asks whether the candidate satisfies the repository's encoded contracts before
it enters the protected branch.

\term{Continuous delivery} automatically produces a verified, releasable
artifact but retains an explicit decision before production publication. The
software is continuously in a state that \emph{could} be released. A reviewer,
scheduled window, regulatory control, or business decision may govern the last
promotion.

\term{Continuous deployment} automatically publishes every qualifying change.
There is no manual approval between the verified candidate and the production
deployment. This does not remove control. Review, tests, branch protection,
scoped permissions, deployment policy, smoke tests, and monitoring remain
controls.

Thus the slash in ``CI/CD'' is ambiguous. State whether \emph{CD} means
delivery or deployment for the system under discussion. This repository uses
both:

\begin{itemize}
  \item the course site uses continuous deployment from reviewed \code{main};
  and
  \item a versioned course release uses continuous delivery to a protected
  approval boundary, followed by an approved publication.
\end{itemize}

Neither choice is intrinsically mature. The appropriate choice follows from
consequence, reversibility, regulation, available evidence, and recovery cost.
A static site can be replaced quickly and carries no database migration. A
clinical model service has a different release boundary.

\section{The running system}

The repository contains four related products:

\begin{enumerate}
  \item LaTeX source under \code{textbook/};
  \item a compiled PDF under \code{output/pdf/};
  \item the \code{rice\_dsm} Python distribution built from \code{src/}; and
  \item a static course site built from reviewed files under \code{site/}.
\end{enumerate}

The textbook is therefore both an explanation of delivery and a delivered
artifact. The workflows are executable course material. Their checked-in YAML
is the authority when an excerpt in the prose and the live configuration ever
differ.

\begin{center}
\begin{tikzpicture}[node distance=7mm]
  \node[flowbox,text width=23mm] (change) {Proposed\\change};
  \node[flowbox,text width=28mm,right=of change] (matrix) {CI matrix\\Linux, macOS, Windows};
  \node[flowbox,text width=23mm,right=of matrix] (gate) {Required\\CI gate};
  \node[flowbox,text width=22mm,right=of gate] (merge) {Review\\and merge};
  \node[flowbox,text width=24mm,below=12mm of matrix] (build) {Build one\\site artifact};
  \node[flowbox,text width=24mm,right=of build] (pages) {Deploy to\\GitHub Pages};
  \node[flowbox,text width=24mm,right=of pages] (smoke) {Verify public\\bytes};
  \draw[flowarrow] (change) -- (matrix);
  \draw[flowarrow] (matrix) -- (gate);
  \draw[flowarrow] (gate) -- (merge);
  \draw[flowarrow] (merge.south) |- (build.north);
  \draw[flowarrow] (build) -- (pages);
  \draw[flowarrow] (pages) -- (smoke);
\end{tikzpicture}
\end{center}

The upper path is integration. The lower path is deployment. Passing the upper
path does not imply that the lower path ran, and a successful upload in the
lower path does not imply that users can retrieve the expected artifact. The
final smoke test addresses that last boundary.

\section{Path one: a pull request becomes verified source}

The \code{course-ci.yml} workflow responds to pull requests, pushes to
\code{main}, and manual dispatch. It creates one job for each supported
operating system. Every job:

\begin{enumerate}
  \item checks out the exact revision without retaining write credentials;
  \item installs the pinned \code{uv} tool and Python version;
  \item reproduces the dependency graph from \code{uv.lock};
  \item prepares the installed package and named notebook kernel;
  \item checks Python style;
  \item builds the wheel and source distribution; and
  \item executes package, repository, and notebook tests.
\end{enumerate}

The matrix is not three redundant green lights. It probes distinct path,
filesystem, shell, process, and environment behavior. A hard-coded POSIX path
may pass on Linux and fail on Windows. A case mismatch may pass on a
case-insensitive filesystem and fail elsewhere.

After the matrix completes, a stable job named \code{CI gate} requires the
entire matrix to succeed. Branch protection can require this one name even when
the internal matrix changes. The gate makes evidence consequential: a failed
candidate cannot enter protected \code{main} through the ordinary merge path.

Other workflows add narrower evidence. The textbook workflow compiles relevant
LaTeX changes, rejects unresolved references and layout warnings, and uploads
the PDF for review. Dependency review examines newly introduced dependency
risk. The pull-request labeler assigns review areas without checking out or
executing untrusted pull-request code.

\begin{checkpointbox}
A pull request receives a green \code{CI gate}, but the new statistical test
uses the wrong null hypothesis. What has CI established? What has it not
established? Identify the missing review evidence without claiming that CI is
useless.
\end{checkpointbox}

\section{Path two: \code{main} becomes a public course site}

The \code{course-pages.yml} workflow is the repository's real continuous
deployment example. A relevant merge to \code{main} triggers the path. Manual
dispatch is also available, but the build job checks that the selected revision
is \code{main}; a feature branch cannot use that route to deploy itself.

\subsection{Build a bounded artifact}

The workflow compiles the textbook, checks the LaTeX log, and runs
\code{scripts/build\_course\_site.py}. The builder creates an empty output tree
containing only:

\begin{itemize}
  \item the generated landing page and fallback page;
  \item reviewed style and icon assets;
  \item the compiled textbook at a stable public path; and
  \item a machine-readable deployment manifest.
\end{itemize}

The builder refuses to merge files into a nonempty output directory. Without
that check, a file left by an earlier run could survive even though the current
source no longer produces it. Clean output is a simple defense against stale
artifact state.

The site artifact is deliberately smaller than the repository. It excludes
tests, source history, local data, caches, notebook outputs, development tools,
and secrets. A deployment boundary should state what enters production rather
than publish whatever happens to be nearby.

\subsection{Record identity in a manifest}

The generated \code{manifest.json} binds the deployment to its inputs and
primary artifact:

\begin{lstlisting}[numbers=none]
{
  "course": "CMOR 438 / INDE 577",
  "package_version": "0.1.0",
  "revision": "8e66ba305f7bb22e5459a4c039999250e29aa35a",
  "generated_at": "2026-08-30T18:00:00Z",
  "artifacts": {
    "textbook": {
      "path": "textbook/data-science-machine-learning-textbook.pdf",
      "bytes": 765614,
      "sha256": "1119dc07f6a9a5cee16444c3272ac67d032c84c4fd1856f77054716595e8691f"
    }
  }
}
\end{lstlisting}

The concrete values above illustrate one local build and will change. The
commit identifier answers which source was built. The byte count catches gross
truncation or substitution. SHA-256 supplies a content digest: changing one
byte is expected to change the digest. A matching digest proves identity with
the recorded bytes under the hash assumption; it does not prove that the book's
mathematics or prose is correct.

\subsection{Promote and then observe}

The build job uploads one Pages artifact. A separate deploy job downloads and
promotes that artifact to the \code{github-pages} environment. It does not
compile the text again. Only that job receives permission to write Pages and
request the short-lived identity token required by the deployment mechanism.

After deployment, \code{scripts/smoke\_test\_course\_site.py} requests the
public page, manifest, and textbook. It checks the page identity, complete
source revision, PDF signature, byte count, and SHA-256 digest. Retries allow a
new edge cache to converge, but retries are bounded. A persistent mismatch is a
failed deployment, not an invitation to loop forever.

An operated model service needs the same identity in its telemetry. The
repository's monitoring laboratory emits a \code{service\_started} structured
event and a \code{rice\_dsm\_model\_info} metric containing model and release
versions. Its dashboard annotates service starts and exposes release identity.
Without this join between deployment and operating evidence, a team cannot
compare symptoms before and after a release or know which artifact produced a
prediction.

This verification distinguishes three claims:

\begin{longtable}{@{}
  >{\raggedright\arraybackslash}p{0.25\textwidth}
  >{\raggedright\arraybackslash}p{0.31\textwidth}
  >{\raggedright\arraybackslash}p{0.34\textwidth}@{}}
\toprule
\textbf{Claim} & \textbf{Evidence} & \textbf{Remaining uncertainty} \\
\midrule
\endhead
The artifact was accepted & Deploy action completed & Public routing or cache
may still be wrong \\
The public site is reachable & HTTP request succeeded & It may expose an older
or incorrect revision \\
The intended PDF is public & Revision, size, and digest match & Content may
still contain an untested error \\
\bottomrule
\end{longtable}

\section{Path three: a version becomes an approved release}

The course site answers ``What is currently published?'' A versioned release
answers a different question: ``Which immutable files constitute version
$v$?'' The \code{course-release.yml} workflow starts only from tags matching
\code{course-v*}.

The version declared in \code{pyproject.toml} is the source of truth. If the
package version is \code{0.1.0}, the release tag must be exactly
\code{course-v0.1.0}. The tag must be annotated rather than lightweight:

\begin{lstlisting}[numbers=none]
git tag -a course-v0.1.0 -m "Course release 0.1.0"
git push origin course-v0.1.0
\end{lstlisting}

An annotated tag is a Git object with a tagger, message, date, and target. The
workflow rejects a mismatched or lightweight tag before publication. It then
reproduces the locked environment, reruns style and test checks, builds the
Python distributions, compiles the textbook, and assembles:

\begin{itemize}
  \item \code{data-science-machine-learning-textbook.pdf};
  \item the versioned \code{rice\_dsm} wheel;
  \item the versioned \code{rice\_dsm} source archive;
  \item \code{release-manifest.json}; and
  \item \code{SHA256SUMS}.
\end{itemize}

The release builder selects exact distribution names for the declared version.
A stale wheel with a similar name in a developer's \code{dist/} directory is
not silently included. The checksum file allows a recipient to recompute the
digest of every release file. Build provenance records that GitHub Actions
produced the subjects from a particular repository workflow and revision.

\begin{professionalbox}
Provenance and checksums answer origin and identity questions. They do not make
malicious source safe, make a weak test suite sufficient, or establish that a
model is scientifically valid. Supply-chain evidence composes with source
review, testing, access control, and domain evidence; it does not replace them.
\end{professionalbox}

\subsection{The approval boundary}

The verified bundle is uploaded once as a workflow artifact. The publication
job names the protected \code{course-release} environment and waits for its
required reviewer. Approval authorizes public release of the already verified
bundle. The job then downloads that exact bundle and attaches its files to the
tag's GitHub Release.

This is continuous delivery because automation reaches a releasable state and
stops at an explicit production decision. The decision should inspect evidence,
not merely confirm that somebody clicked a button. A useful approval view
includes the source revision, tag and project version, completed checks,
artifact names and digests, known limitations, and rollback plan.

If publication is interrupted after creating the GitHub Release but before all
attachments arrive, a rerun replaces the attachments with files from the
verified workflow artifact. Idempotent or repairable operations matter because
distributed systems can fail between two externally visible steps.

\section{Build once, promote the same artifact}

Suppose a build job tests PDF $A$, then an approval job checks out the same
commit and recompiles PDF $B$ for publication. Even with identical source, the
outputs may differ because a dependency, compiler, clock, generated input, or
configuration changed. Approval evidence about $A$ does not necessarily apply
to $B$.

The safer rule is:

\begin{quote}
Build once. Identify the artifact. Verify it. Promote those same bytes.
\end{quote}

Configuration may vary by deployment environment, but changing configuration
should not silently mutate the artifact. A container image digest, wheel hash,
or PDF checksum provides a stable subject for evidence. This separation also
makes rollback concrete: retain a known-good artifact and its compatible
configuration rather than hoping an old source revision can be rebuilt later.

Reproducible builds strengthen the design further by allowing independent
builds to produce the same bytes. They do not eliminate the need to promote the
verified artifact; reproducibility itself must be measured and maintained.

\section{Authority is part of the pipeline}

Automation that can only read source cannot publish. Automation that can write
every repository resource has unnecessary authority. The workflows declare a
read-only default and grant additional permissions to the narrow job that needs
them.

\begin{longtable}{@{}
  >{\raggedright\arraybackslash}p{0.25\textwidth}
  >{\raggedright\arraybackslash}p{0.26\textwidth}
  >{\raggedright\arraybackslash}p{0.39\textwidth}@{}}
\toprule
\textbf{Job} & \textbf{Additional authority} & \textbf{Reason} \\
\midrule
\endhead
CI and build jobs & None beyond reading source & Untrusted or fallible build
steps should not publish \\
Pages deploy job & Pages write and short-lived identity token & Promote the one
uploaded Pages artifact \\
Release provenance job & Attestation write and identity token & Bind provenance
to built subjects \\
Release publication job & Repository contents write & Create or repair assets
on the tagged GitHub Release \\
\bottomrule
\end{longtable}

The repository does not need a personal access token, AWS key, or other
long-lived credential for these paths. GitHub supplies a short-lived repository
token to each job, and job-level permissions narrow its use. Third-party actions
still execute code with the job's available authority, so actions are pinned to
immutable commit identifiers and updated through review.

If a later exercise deploys to a cloud provider, prefer workload identity or an
equivalent short-lived exchange. Store unavoidable secrets in the platform's
secret facility, scope them to the environment, prevent them from reaching
untrusted pull-request code, mask them in logs, and plan rotation before an
incident.

\section{Failure is a state to diagnose}

A pipeline is a state machine with evidence at each transition. ``CI/CD
failed'' is too imprecise to diagnose.

\begin{longtable}{@{}
  >{\raggedright\arraybackslash}p{0.23\textwidth}
  >{\raggedright\arraybackslash}p{0.31\textwidth}
  >{\raggedright\arraybackslash}p{0.36\textwidth}@{}}
\toprule
\textbf{Observed state} & \textbf{Likely boundary} & \textbf{First evidence} \\
\midrule
\endhead
One matrix job fails & Platform-specific source or environment behavior & First
causal traceback on that runner \\
All tests pass; site build fails & Artifact construction contract & Missing
input, unresolved template marker, nonempty output \\
Deploy action fails & Hosting authority or environment policy & Deployment job,
token permission, Pages configuration \\
Deploy succeeds; smoke fails & Public routing, cache, or artifact mismatch &
Requested URL, manifest revision, byte count, digest \\
Release build rejects tag & Version identity policy & Tag object type and
\code{pyproject.toml} version \\
Release waits & Protected environment decision & Reviewer and environment
policy \\
Release exists with incomplete files & Interrupted publication & Release assets
versus verified workflow artifact \\
\bottomrule
\end{longtable}

Start with the earliest causal failure, not the last cascade message. Reproduce
locally when the failing boundary exists locally. A Pages permission error
cannot be reproduced by rerunning a unit test; a malformed manifest can.

Flaky tests deserve diagnosis. Repeatedly rerunning until green changes the
sample, hides uncertainty, and teaches the gate to accept noise. Record the
failure, isolate nondeterminism, repair the test or system, and preserve a
regression case.

\section{Rollback and roll-forward}

Static-site rollback is unusually simple because the artifact has no mutable
database state. The normal repair is a focused revert or corrective pull
request, followed by the same CI and deployment path used for forward changes.
This preserves review, evidence, and an auditable transition. Directly editing
files on a server creates unreviewed state that the next deployment will erase.

A published version tag should not be deleted and moved to different source.
Doing so makes two consumers receive different bytes under the same name. Fix a
bad \code{course-v0.1.0} release by incrementing the project to
\code{0.1.1}, merging the correction, and publishing
\code{course-v0.1.1}. This is a roll-forward.

Rollback becomes harder when a deployment mutates durable state. Old and new
application versions may overlap during a rolling release. A database migration
must often be backward-compatible until both versions no longer run. A model
rollback may fail if the feature schema or population has changed. Recovery
design must therefore name the state being reversed, compatibility assumptions,
owner, trigger, and evidence of rehearsal.

\section{A complete worked change}

Consider a pull request that corrects one equation and changes one reusable
Python function used by a notebook.

\begin{enumerate}
  \item \textbf{Propose.} The branch records the source diff. The pull request
  explains the mathematical correction, interface consequence, and relevant
  tests.
  \item \textbf{Integrate.} The CI matrix rebuilds the package and executes
  every notebook. The textbook workflow compiles the PDF and makes it available
  for visual review. Neither artifact is public production state.
  \item \textbf{Review.} A reviewer examines the equation, API behavior, test
  evidence, and rendered page. The required \code{CI gate} passes.
  \item \textbf{Merge.} The protected branch advances to the reviewed commit.
  This event makes the revision eligible for site deployment.
  \item \textbf{Build.} The Pages workflow compiles the text and creates one
  bounded site artifact with a manifest naming the merge commit and PDF digest.
  \item \textbf{Deploy.} The Pages job promotes the artifact to the public
  environment with narrowly scoped authority.
  \item \textbf{Observe.} The smoke job requests the public artifacts and
  compares their identities with the manifest and expected commit.
  \item \textbf{Release when intended.} If the change belongs to a named course
  version, the maintainer increments \code{pyproject.toml} through a pull
  request, creates the matching annotated tag after merge, inspects the release
  evidence, and approves publication.
\end{enumerate}

The example contains two distinct public transitions. The site changes after
merge; the immutable version changes only after a tag and approval. A team can
choose other policies, but it must not blur these identities.

\section{Inspect the system locally}

Local execution cannot reproduce GitHub's identity tokens, environment
approval, or Pages infrastructure. It can reproduce artifact construction and
verification logic.

Build the text and site from the repository root:

\begin{lstlisting}[numbers=none]
make -C textbook
uv run python scripts/build_course_site.py \
  --revision local-preview \
  --textbook output/pdf/data-science-machine-learning-textbook.pdf
uv run python -m http.server 8000 --directory build/course-site
\end{lstlisting}

Open \url{http://localhost:8000} and inspect the page and
\code{manifest.json}. Stop the server with \code{Ctrl+C}. Before rebuilding,
remove only the generated \code{build/course-site/} directory; the builder's
refusal to mix old and new output is intentional.

The following investigation distinguishes mechanisms:

\begin{enumerate}
  \item Recompute the PDF SHA-256 digest with an operating-system tool or a
  short Python program and compare it with the manifest.
  \item Change a copy of one byte in the PDF and recompute the digest. Do not
  modify the tracked course PDF for this experiment.
  \item Explain why a new digest detects the change but cannot say whether the
  change improved the book.
  \item Identify which workflow-only controls were absent from the local run.
\end{enumerate}

\section{Design exercises}

\begin{enumerate}
  \item Draw the state machine for the Pages workflow. Label every transition
  with its triggering event, required evidence, and authority.
  \item Suppose the public page identifies the expected commit but the PDF hash
  differs. Give three non-equivalent hypotheses and the observation that would
  distinguish each.
  \item Design a staging environment for a model API. State which artifact is
  promoted, which configuration varies, which tests run after deployment, and
  which evidence permits production release.
  \item A release approver sees green tests but no visual PDF review. Should the
  release proceed? State the claim the missing review protects.
  \item Extend the failure table for a database migration that succeeds on one
  replica and fails on another. Identify why restoring application code alone
  may be insufficient.
  \item Compare a manual production approval with a fully automatic policy.
  Name one system for which each is reasonable and justify the difference from
  consequence and recovery evidence.
  \item Write a release dossier for \code{course-v0.1.0}: source revision,
  artifacts, digests, checks, approver, limitations, and recovery action.
\end{enumerate}

\section{Operational checklist}

Before calling a pipeline complete, ask:

\begin{itemize}
  \item Is the source revision immutable and recorded?
  \item Does the build start from a clean, locked environment?
  \item Is each artifact bounded, named, and content-addressable by a digest?
  \item Do later stages promote the verified bytes without rebuilding?
  \item Are tests attached to risks at unit, integration, contract, system, and
  deployed boundaries as appropriate?
  \item Can untrusted source obtain a secret or publication permission?
  \item Does each job have only the authority required for its stage?
  \item Is the release decision automatic or manual, and is that policy stated?
  \item Does post-deployment verification observe the public boundary?
  \item Is recovery possible with a known-good artifact and compatible state?
  \item Will a future investigator be able to connect public behavior to source,
  configuration, evidence, and decision?
\end{itemize}

\section{Summary}

CI tests an integrated source candidate. Continuous delivery produces a
verified releasable artifact and retains an explicit production decision.
Continuous deployment publishes every qualifying change automatically. This
repository uses all three ideas: a cross-platform required gate, automatic
course-site deployment from protected \code{main}, and an approved tagged
release of the textbook and Python package.

The central engineering invariant is not a product name or YAML keyword. It is
the evidence chain:

\begin{quote}
immutable source $\longrightarrow$ bounded artifact $\longrightarrow$
verification $\longrightarrow$ authorized promotion $\longrightarrow$
operational observation.
\end{quote}

Breaking any arrow weakens the claim that a public result is the one that was
reviewed and tested.
